% The part trap, made concrete in one register. Same field, same bits, same
% meaning, different symbol.
% \bitfield always draws at y = 0, so each register lives in its own scope, and
% a two-bit field is 8 mm wide, so its name is called out below rather than
% written inside it.
\begin{tikzpicture}[font=\sffamily\small]

% ================================================================ this part
\begin{scope}
  \node[regname] at (-2mm,3mm) {this part};
  \bitscale
  \bitfield{31}{30}{regrsvd}{}
  \bitfield{29}{28}{regset}{}
  \bitfield{27}{0}{regcell}{other kernel clock selectors, not shown}
  \draw[busstub] (12mm,-0.5mm) -- (12mm,-5mm);
  \node[font=\sffamily\scriptsize, anchor=north] at (12mm,-5mm) {FDCANSEL};
  \node[font=\tiny\ttfamily, anchor=west] at (24mm,-6.5mm) {RCC\_CDCCIP1R, bits 29:28};
  \node[chlabel, anchor=west, text width=34mm] at (132mm,0mm)
    {named after this part's two power domains};
\end{scope}

% ================================================================ the sibling
\begin{scope}[yshift=-30mm]
  \node[regname] at (-2mm,3mm) {the sibling};
  \bitscale
  \bitfield{31}{30}{regrsvd}{}
  \bitfield{29}{28}{regset}{}
  \bitfield{27}{0}{regcell}{other kernel clock selectors, not shown}
  \draw[busstub] (12mm,-0.5mm) -- (12mm,-5mm);
  \node[font=\sffamily\scriptsize, anchor=north] at (12mm,-5mm) {FDCANSEL};
  \node[font=\tiny\ttfamily, anchor=west] at (24mm,-6.5mm) {RCC\_D2CCIP1R, bits 29:28};
  \node[chlabel, anchor=west, text width=34mm] at (132mm,0mm)
    {named after its three power domains, which this part does not have};
\end{scope}

% ================================================================ the reading
\node[umlnote, text width=62mm, anchor=north west] at (0mm,-46mm)
  {\textbf{The same two bits, at the same positions, with the same meaning.} A
   line copied from material written for the sibling part does not silently do
   the wrong thing here: it fails to compile, because the symbol does not exist.
   That is the best failure mode available, and it is worth knowing about before
   an afternoon is spent on it.};

\node[umlnote, text width=64mm, anchor=north west] at (66mm,-46mm)
  {\textbf{The wider correction this volume carries.} The warning that this part
   is not the family member most material was written for is true at board,
   clock and clock-register level. At the bus peripheral itself it is false: a
   comparison of the vendor's own device headers shows the same licensed
   controller core, at the same base addresses, with the same calibration unit,
   the same message memory base and identical bit-timing field positions.};

\node[chlabel, anchor=north west, text width=130mm] at (0mm,-86mm)
  {So the bus material written for the sibling is directly reusable, and three
   things are not: this register's name, the domain naming around it, and the
   maximum core frequency that changes the clock arithmetic. A reader can repeat
   the comparison in ten minutes, because both headers are public, permissively
   licensed and fetchable, which is why this volume settles silicon questions
   from them rather than from prose.};
\end{tikzpicture}
